\subsection{Introduction to the FJ model}
\label{subsec:fj-intro}

\subsubsection{The FJ model}
\label{subsubsec:fj-model}

In the FDG model every agent keeps averaging the opinions it listens to, and a
persistent disagreement can only come from the structure of the network, i.e.\ from
several closed groups (Theorem~\ref{thm:disagreement}). The Friedkin--Johnsen (FJ)
model adds a second ingredient: each agent can also be tied to a fixed opinion of its own,
called its \emph{prejudice}. Following~\cite{proskurnikov2017tutorial}, the model reads
\begin{equation}\label{eq:fj}
  x(k+1) = \Lambda W x(k) + (I - \Lambda)\, u ,
\end{equation}
where
\begin{itemize}
  \item $W \in \mathbb{R}^{n \times n}$ is a row-stochastic matrix, the \emph{matrix of
        social influences}, i.e.\ the weighted adjacency matrix of the influence
        digraph $G$ (Definition~\ref{def:adjacency}), which plays the role of $A$ in
        the FDG model;
% REVISED (2026-10-28, review): the edge convention is already stated in Section 1.1. Was:
%   \item $W \in \mathbb{R}^{n \times n}$ is a row-stochastic matrix, the \emph{matrix of
%         social influences}, i.e.\ the weighted adjacency matrix of the influence
%         digraph $G$ (Definition~\ref{def:adjacency}); with the convention of
%         Section~\ref{subsec:prelim},
%         $w_{ij} > 0$ means that agent $i$ is influenced by agent $j$, so $W$ plays the
%         role that $A$ plays in the FDG model;
  \item $\Lambda = \operatorname{diag}(\lambda_1, \dots, \lambda_n)$, with
        $\lambda_i \in [0,1]$ the \emph{susceptibility} of agent $i$ to social
        influence;
  \item $u \in \mathbb{R}^n$ is a constant vector collecting the agents'
        \emph{prejudices}.
\end{itemize}
Componentwise,
\begin{equation}\label{eq:fj-i}
  x_i(k+1) = \lambda_i \sum_{j=1}^{n} w_{ij}\, x_j(k) + (1 - \lambda_i)\, u_i .
\end{equation}
At every step agent $i$ averages the opinions it listens to, exactly as in the FDG
model, and mixes this average with its prejudice: the weight $\lambda_i$ goes to the
others and the complement $1 - \lambda_i$ goes to $u_i$. The two quantities play
different roles: the prejudice $u_i$ is an opinion and says \emph{where} the agent is
anchored, while $1 - \lambda_i$ is a weight and says \emph{how strongly}. The model with
$u = x(0)$ appears in~\cite{bullo2024lns}, where $\lambda_i$ is called the openness
level of agent $i$ and $1 - \lambda_i$ its attachment to its initial opinion, which
plays the role of the prejudice.
% REVISED (2026-10-28, review): CORRECTION: the model of the book (Exercise E5.24) has u = x(0), it is not the same model with an independent u. Was:
% The same
% model appears in~\cite{bullo2024lns}, where $\lambda_i$ is called the openness level of
% agent $i$ and $1 - \lambda_i$ its attachment to the prejudice.

The special cases of~\eqref{eq:fj} are the following~\cite{proskurnikov2017tutorial}.
\begin{itemize}
  \item If $\lambda_i = 1$, agent $i$ follows the FDG rule. If $\Lambda = I$, the FJ
        model reduces to the FDG model with $A = W$; for this reason we assume
        $\Lambda \neq I$ from now on.
  \item If $\lambda_i < 1$, agent $i$ is called \emph{prejudiced}: it is anchored to its
        prejudice $u_i$.
  \item If $\lambda_i = 0$, then $x_i(k) = u_i$ for all $k \geq 1$: the agent ignores
        the others and jumps to its prejudice after one step. Its opinion never changes,
        i.e.\ the agent is stubborn in the sense of Definition~\ref{def:stubborn}, only
        if in addition $u_i = x_i(0)$. A stubborn agent of the FJ model is therefore an
        agent with $\lambda_i = 0$ and $u_i = x_i(0)$.
\end{itemize}

In social terms, the self-weight of the FDG model and the prejudice of the FJ model are
two different forms of resistance to influence. A self-weight ties an agent to its
\emph{current} opinion, which still moves towards the others
(Section~\ref{subsubsec:fdg-stub-model}); a prejudice ties it to a \emph{fixed}
opinion, which never moves.

In the original formulation of Friedkin and Johnsen the prejudices coincide with the
initial opinions, $u = x(0)$, and the susceptibilities are coupled to the self-weights
through $1 - \lambda_i = w_{ii}$~\cite{proskurnikov2017tutorial,bullo2024lns}. As
in~\cite{proskurnikov2017tutorial}, we do not impose these conditions: $u$ and $x(0)$
are independent, and so are $\Lambda$ and $W$. The choice $u = x(0)$ will be needed in
Section~\ref{subsubsec:fj-centrality} to define the influence centrality.

\subsubsection{Convergence and stability}
\label{subsubsec:fj-conv}
% RESTRUCTURED (2026-10-23): this subsubsection was called "Convergence: all agents
% P-dependent" and started from Corollary 22 of the tutorial. It now follows the order
% of the tutorial: definitions, P-dependent agents and decomposition, Theorem 21 (where
% V is defined), Corollary 22 (all agents P-dependent), Corollary 23, and only at the
% end our observations (error dynamics). The former subsubsection "The general case:
% P-independent agents" (label subsubsec:fj-general) is merged here.
% REVISED (2026-10-24): Definition P-dependent, Theorem 21, Corollaries 22 and 23 are
% now stated exactly as in the tutorial (only the notation is adapted); x(inf) = Vu is
% now outside Corollary 22, in eq. fj-V-all.
% REVISED (2026-10-23, after the general brainstorming): the observations are no longer
% in a separate block at the end; the intuition for Corollary 22 follows it directly,
% and Corollary 23 closes the subsubsection, as in the tutorial. The duplicated formula
% of V (former eq. fj-V-all before the corollary) is now only in the corollary. The
% version before this revision is saved in fj/fj_intro_OLD_2026-10-23.tex.

The FJ model~\eqref{eq:fj} is a linear system with the constant input $u$. Following
\cite{proskurnikov2017tutorial}, we say that it \emph{converges} if the limit
$x(\infty) = \lim_{k \to +\infty} x(k)$ exists for every initial condition $x(0)$, and
that it is \emph{asymptotically stable} if $\Lambda W$ is convergent, i.e.\
$\rho(\Lambda W) < 1$ (Definition~\ref{def:semiconv} and
Theorem~\ref{thm:spectral-radius}). Regular matrices are recalled in Definition~\ref{def:regular}.
% REVISED (2026-10-24): the prose definition of "regular" ("Following the same reference,
% a matrix A is called regular if ...") is replaced by Definition def:regular in the
% Preliminaries (Definition 15 of the tutorial).

Both properties are characterized by distinguishing the agents that are reached by the
prejudices from those that are not.

\begin{definition}[P-dependent agents~\cite{proskurnikov2017tutorial}]
\label{def:p-dependent}
Agent $i$ is \emph{P-dependent} if it is prejudiced or influenced by some prejudiced
agent $j$, that is, a walk from $i$ to $j$ in the graph $G$ exists. Otherwise, the
agent is \emph{P-independent}.
\end{definition}

The walk goes from the agent to the prejudiced agent because of our convention on the
edges; in~\cite{proskurnikov2017tutorial}, where the edges are drawn in the opposite
direction, the same walk goes from the prejudiced agent to the agent. In social terms,
an agent is P-dependent if at least one prejudice reaches it.

A P-independent agent $i$ is not prejudiced, so $\lambda_i = 1$, and it is not
influenced by any P-dependent agent: if $w_{ij} > 0$ with $j$ P-dependent, there would
be a walk from $i$ to a prejudiced agent through $j$, and $i$ would be P-dependent.
Since $\Lambda \neq I$, there is at least one prejudiced agent, hence at least one
P-dependent agent. Renumbering the agents so that agents $1, \dots, r$, with
$r \geq 1$, are P-dependent and agents $r+1, \dots, n$ are P-independent, and
splitting $x = [x^1; x^2]$, $u = [u^1; u^2]$, $W$ and $\Lambda$ accordingly, we
have $W^{21} = 0$ and $\Lambda^{22} = I$, and the model~\eqref{eq:fj} decomposes
as~\cite{proskurnikov2017tutorial}
\begin{align}
  x^1(k+1) &= \Lambda^{11}\bigl[W^{11} x^1(k) + W^{12} x^2(k)\bigr]
              + (I - \Lambda^{11})\, u^1, \label{eq:fj-dep}\\
  x^2(k+1) &= W^{22} x^2(k), \label{eq:fj-indep}
\end{align}
where $W^{22}$ is row-stochastic, since $W$ is row-stochastic and $W^{21} = 0$. The
P-independent agents therefore follow an FDG
model among themselves, with $A = W^{22}$, and no prejudice reaches them; the
P-dependent agents receive the opinions of the P-independent ones as an additional
input.

\begin{theorem}[\cite{proskurnikov2017tutorial}]
\label{thm:fj-general}
Let the community have $r \geq 1$ P-dependent agents and $n - r \geq 0$
P-independent ones. Then, the subsystem~\eqref{eq:fj-dep} is asymptotically stable,
i.e.\ $\Lambda^{11} W^{11}$ is Schur stable $\rho(\Lambda^{11} W^{11}) < 1$. The
model~\eqref{eq:fj} is convergent if and only if $r = n$ or~\eqref{eq:fj-indep} is
convergent, i.e.\ $W^{22}$ is regular. In this case
\begin{equation}\label{eq:fj-V}
  x^1(\infty) = V \begin{bmatrix} u^1 \\ x^2(\infty) \end{bmatrix}, \qquad
  V \triangleq (I_r - \Lambda^{11} W^{11})^{-1}
      \begin{bmatrix} I_r - \Lambda^{11} & \Lambda^{11} W^{12} \end{bmatrix}.
\end{equation}
The matrix $V$ is stochastic, i.e.\ the final opinion of any agent is a convex
combination of the prejudices and final opinions of P-independent agents.
\end{theorem}

In~\cite{proskurnikov2017tutorial}, \emph{stochastic} stands for row-stochastic in the
sense of Definition~\ref{def:row-stochastic}.

Theorem~\ref{thm:fj-general} can be read according to the number $r$ of P-dependent
agents.
\begin{itemize}
  \item $r = 0$. This case is not covered by Theorem~\ref{thm:fj-general}: no agent is
        prejudiced, i.e.\ $\Lambda = I$, and the FJ model~\eqref{eq:fj} reduces to the
        FDG model with $A = W$ (Section~\ref{subsec:fdg-intro}). This is why we assume
        $\Lambda \neq I$, which guarantees $r \geq 1$.
  \item $0 < r < n$. P-dependent and P-independent agents coexist. The P-independent
        agents follow~\eqref{eq:fj-indep}, which is an FDG model with the
        row-stochastic matrix $W^{22}$, so the outcomes of Table~\ref{tab:rosetta}
        apply to them. Denoting by $C^{22}$ the condensation digraph of the subgraph of
        P-independent agents, whose adjacency matrix is $W^{22}$, they reach
        \begin{itemize}
          \item a consensus, if $C^{22}$ has a unique sink, which is aperiodic
                (cases~(i) and~(iii));
          \item a disagreement, if $C^{22}$ has several sinks, all aperiodic
                (case~(iv));
          \item no convergence, if at least one sink of $C^{22}$ is periodic
                (cases~(ii) and~(v)).
        \end{itemize}
        In the first two cases $W^{22}$ is regular, the model~\eqref{eq:fj} converges
        and the final opinions of the P-dependent agents are given
        by~\eqref{eq:fj-V}; in the third case the model~\eqref{eq:fj} does not
        converge.
  \item $r = n$. There are no P-independent agents, the subsystem~\eqref{eq:fj-dep} is
        the whole model and the model~\eqref{eq:fj} always converges. This case is the
        subject of Corollary~\ref{cor:fj-stable} below.
\end{itemize}
In social terms, the P-dependent agents settle on a compromise between the prejudices
and the opinions of the P-independent agents they listen to whenever these opinions
settle; whether the whole group settles is decided only by the P-independent agents.
% REVISED (2026-10-24): reading of Theorem 21 organized by the number r of P-dependent
% agents (r = 0: FDG, excluded; 0 < r < n: outcomes of Table 1 for the P-independent
% agents; r = n: Corollary 22). The symbol C^{22} is ours.


\begin{corollary}[\cite{proskurnikov2017tutorial}]
\label{cor:fj-stable}
The Friedkin--Johnsen model~\eqref{eq:fj} is asymptotically stable if and only if all
agents are P-dependent. Then, $V = (I - \Lambda W)^{-1}(I - \Lambda)$. This holds, in
particular, if $\Lambda < I_n$ or $\Lambda \neq I_n$ and $G$ is strongly connected
($W$ is irreducible).
\end{corollary}

With $r = n$,~\eqref{eq:fj-V} gives, for every $x(0)$,
\begin{equation}\label{eq:fj-V-all}
  x(\infty) = V u ,
\end{equation}
with $V$ as in Corollary~\ref{cor:fj-stable}.

The condition of Corollary~\ref{cor:fj-stable} can also be read on the matrix
$\Lambda W$ with the tools of the course. The matrix $\Lambda W$ is no longer
row-stochastic: since $W$ is row-stochastic, the $i$-th row of $\Lambda W$ sums to
$\lambda_i$. Rows of agents that are not prejudiced sum to $1$, rows of prejudiced
agents sum to less than $1$, and since $\Lambda \neq I$ at least one such row exists.
Hence $\Lambda W$ is row-substochastic (Definition~\ref{def:row-substochastic}) and
Lemma~\ref{lem:substochastic} applies.
% MOVED (2026-10-23): Definition "Row-substochastic matrix" and Lemma "Convergent
% row-substochastic matrices" were here; they are now in the Preliminaries
% (Section 1.1.2), with the definition of weighted out-degree.

For the matrix $\Lambda W$, the weighted out-degree of agent $i$ is $\lambda_i$: the nodes with
out-degree $1$ are the agents that are not prejudiced, and the nodes with out-degree
less than $1$ are the prejudiced agents. The condition of
Lemma~\ref{lem:substochastic} therefore reads: every agent that is not prejudiced is
influenced, directly or through intermediaries, by a prejudiced agent, i.e.\ all the
agents are P-dependent. This is exactly the condition of
Corollary~\ref{cor:fj-stable}. Indeed,~\cite{proskurnikov2017tutorial} remarks that
Corollary~\ref{cor:fj-stable} can be turned into a criterion of Schur stability for
substochastic matrices, and~\cite{bullo2024lns} proves the same condition for the FJ
model through this lemma.

Two observations of ours give an intuition for Corollary~\ref{cor:fj-stable}. First,
letting $k \to +\infty$ in~\eqref{eq:fj} shows that the limit $x(\infty) = Vu$ is an
equilibrium, $x(\infty) = \Lambda W x(\infty) + (I - \Lambda)u$. Letting
$e(k) = x(k) - x(\infty)$ and subtracting this equation from~\eqref{eq:fj} gives
\begin{equation}\label{eq:fj-error}
  e(k+1) = \Lambda W\, e(k),
\end{equation}
an autonomous linear system in which the prejudice $u$ does not appear: $\Lambda W$
decides whether the opinions converge, $u$ decides where.

Second, the linear system~\eqref{eq:fj-error} shows the role of the P-dependent
agents. At each step, agent $i$ keeps a fraction $\lambda_i$ of the weighted average
of the deviations of the agents it listens to. An agent that is not prejudiced
($\lambda_i = 1$) only averages the deviations, while a prejudiced agent loses a
fraction $1-\lambda_i$ of them (in~\eqref{eq:fj}, this fraction is replaced by the
agent's fixed prejudice). The deviations therefore leave the group only through the
prejudiced agents, and they vanish only if every agent is influenced, directly or
through intermediaries, by one of them, i.e.\ all agents are P-dependent. A group of
agents that is not influenced, even indirectly, by any prejudiced agent loses no
deviation and behaves as an FDG model, as in~\eqref{eq:fj-indep}.

In social terms, the group converges, for every initial opinion, as soon as every agent
that is not anchored to a prejudice listens, directly or through intermediaries, to
someone who is. On a strongly connected network, such as the two networks of
Figure~\ref{fig:networks}, a single prejudiced agent is enough. Note that this does not
require aperiodicity, unlike the FDG model of Section~\ref{subsec:fdg-stub}.

If some agents are P-independent, the model is not asymptotically stable, but it may
still converge. Theorem~\ref{thm:fj-general} gives the condition in terms
of $W^{22}$; in terms of the whole matrix $\Lambda W$ it reads as follows.

\begin{corollary}[\cite{proskurnikov2017tutorial}]
\label{cor:fj-regular}
The model~\eqref{eq:fj} converges if and only if $\Lambda W$ is regular, i.e.\ the
limit $\lim_{k \to \infty} (\Lambda W)^k$ exists.
\end{corollary}

As remarked in~\cite{proskurnikov2017tutorial}, this property is not trivial: a system
$x(k+1) = A x(k) + B u$ with a regular matrix $A$ may not converge, and may even have
unbounded solutions, as for $A = B = I$.

\subsubsection{How the opinions converge: the matrix \texorpdfstring{$V$}{V} and the augmented graph}
\label{subsubsec:fj-final}
% RENAMED (2026-10-24): was "Final opinions: consensus or disagreement" -- the tutorial
% does not use consensus/disagreement for the FJ model (Section 6); consensus
% (Definition 14 of the tutorial) is a property of the FDG model.
% RENAMED (2026-10-23): was "The equilibrium and the augmented graph"
% (label subsubsec:fj-equilibrium).

When all agents are P-dependent, the final opinions are given
by~\eqref{eq:fj-V-all}, and the convex combination of Theorem~\ref{thm:fj-general}
involves only the prejudices:
\[
  x_i(\infty) = \sum_{j=1}^{n} v_{ij}\, u_j .
\]
Three facts follow.
\begin{itemize}
  \item The initial opinions are forgotten: $x(0)$ does not appear in $x(\infty)$.
  \item Only the prejudices of prejudiced agents count: if $\lambda_j = 1$, the $j$-th
        column of $I - \Lambda$ is zero, hence so is the $j$-th column of $V$, and
        $u_j$ has no weight in any final opinion.
  \item Every final opinion lies between the smallest and the largest prejudice of the
        prejudiced agents.
\end{itemize}
The comparison with the FDG model is instructive. There, the consensus value
$w^\top x(0)$ is a weighted average of the initial opinions, with the same weights $w$
for every agent (Theorem~\ref{thm:consensus}); here, the final opinions are weighted
averages of the prejudices, and each agent has its own weights, given by its row of
$V$.

In the FDG model, consensus is a property of the model: the model reaches a consensus
if the final opinions coincide for any initial condition
$x(0)$~\cite{proskurnikov2017tutorial}. In the FJ model with all agents P-dependent,
the limit $x(\infty) = Vu$ does not depend on $x(0)$
(equation~\eqref{eq:fj-V-all}): whether the final opinions coincide is decided by the
prejudices $u$, which are an external input, through the matrix $V$, and no longer by
the initial opinions. In our
reading, this is why~\cite{proskurnikov2017tutorial} studies the FJ model in terms of convergence and
of the form of the limit, and not in terms of consensus.
% REVISED (2026-10-24): was "...decided by the prejudices u, which are an external
% input, and not by the network" -- false in general: with a single prejudiced agent
% V = 1 e_1^T (V row-stochastic, zero columns for lambda_j = 1) and the final opinions
% coincide for every u (Section 3.2.2), so the network matters too.

The notion of consensus of the FDG model applies again only to the P-independent agents
($0 < r < n$), which evolve as an FDG model with matrix $W^{22}$ and whose limit depends
on their initial opinions $x^2(0)$ (Section~\ref{subsubsec:fj-conv}). The P-dependent
agents converge to convex combinations of the prejudices and of the final opinions of
the P-independent agents (Theorem~\ref{thm:fj-general}), so they do not necessarily end
up at the consensus value of the P-independent agents, even when there is one.
% REVISED (2026-10-28, review): shortened: the three outcomes for the P-independent agents are already listed in Section 3.1.2. Was:
% The notion of consensus of the FDG model applies again only when some agents are
% P-independent ($0 < r < n$), and only to them. These agents are not prejudiced and are
% not influenced by any P-dependent agent, so they evolve as an FDG model with matrix
% $W^{22}$ (equation~\eqref{eq:fj-indep}), and their limit depends on their initial
% opinions $x^2(0)$. For them, the three outcomes of the FDG model are possible:
% consensus, disagreement or no convergence (Table~\ref{tab:rosetta}). The P-dependent
% agents, instead, converge to convex combinations of the prejudices and of the final
% opinions of the P-independent agents (Theorem~\ref{thm:fj-general}); hence, even if
% the P-independent agents reach a consensus among themselves, the P-dependent agents do
% not necessarily end up at the same value.

The limit~\eqref{eq:fj-V-all} has a direct interpretation in terms of the FDG model.
Following~\cite{proskurnikov2017tutorial}, the FJ model~\eqref{eq:fj} may be
considered as an augmented FDG model with $n$ virtual stubborn agents, anchored at the
prejudices: the virtual agent $i'$ associated with agent $i$ keeps the opinion
$x_{i'}(k) = u_i$ for all $k$. Explicitly (this rewriting is ours), stacking the
opinions $x$ of the real agents and $x'$ of the virtual agents,
\begin{equation}\label{eq:fj-augmented}
  \begin{bmatrix} x(k+1) \\ x'(k+1) \end{bmatrix}
  = \begin{bmatrix} \Lambda W & I - \Lambda \\ 0 & I \end{bmatrix}
    \begin{bmatrix} x(k) \\ x'(k) \end{bmatrix},
  \qquad x'(0) = u .
\end{equation}
Row $i$ of the matrix in~\eqref{eq:fj-augmented} sums to $\lambda_i + (1-\lambda_i) = 1$,
and row $i'$ sums to $1$; hence the matrix is row-stochastic
and~\eqref{eq:fj-augmented} is an FDG model on $2n$ agents. In its digraph, which we
call the \emph{augmented graph}:
\begin{itemize}
  \item a real agent $i$ is influenced by a real agent $j$ with weight
        $\lambda_i w_{ij}$;
  \item a real agent $i$ is influenced by its virtual agent $i'$ with weight
        $1 - \lambda_i$, so this edge exists only if agent $i$ is prejudiced;
  \item a virtual agent $i'$ is influenced only by itself, with weight $1$, so its
        opinion never changes: it is stubborn (Definition~\ref{def:stubborn}).
\end{itemize}
Figure~\ref{fig:augmented} shows the FJ model on the directed cycle of
Figure~\ref{fig:networks}(a) when agent~$1$ is the only prejudiced agent, and the
corresponding augmented graph. In panel~(a) the prejudice is drawn as an input
entering agent~$1$, since it is not an agent of the network but an external piece of
information; in panel~(b) it becomes the virtual agent $1'$, and the edge follows our
convention: it goes from $1$ to $1'$, because agent~$1$ listens to $1'$. In social
terms, a prejudice behaves as an additional agent that never changes its opinion: a
prejudiced agent keeps listening to the others, and also listens to this stubborn
agent placed beside it.

\begin{figure}[H]
  \centering
  \begin{tikzpicture}[baseline=(c.base)]
    \foreach \i in {1,...,6} {
      \node[agent] (p\i) at ({90-60*(\i-1)}:1.3) {\i};
    }
    \node (c) at (0,0) {};
    \foreach \i/\j in {2/3,3/4,4/5,5/6,6/1} {
      \draw[edge] (p\i) -- (p\j);
    }
    \draw[edge] (p1) -- node[weight, above right] {$\lambda_1$} (p2);
    \draw[edge] (0,2.7) node[above, font=\scriptsize] {$u_1$}
      -- node[weight, right] {$1-\lambda_1$} (p1);
    \node at (0,-2.0) {\small (a) FJ model on the directed cycle};
  \end{tikzpicture}
  \hspace{1.5cm}
  \begin{tikzpicture}[baseline=(c.base)]
    \foreach \i in {1,...,6} {
      \node[agent] (a\i) at ({90-60*(\i-1)}:1.3) {\i};
    }
    \node (c) at (0,0) {};
    \foreach \i/\j in {2/3,3/4,4/5,5/6,6/1} {
      \draw[edge] (a\i) -- (a\j);
    }
    \draw[edge] (a1) -- node[weight, above right] {$\lambda_1$} (a2);
    \node[agent, dashed, minimum size=5.5mm] (v1) at (0,2.7) {$1'$};
    \draw[edge] (a1) -- node[weight, right] {$1-\lambda_1$} (v1);
    \draw[selfloop] (v1) to node[weight, above] {$1$} (v1);
    \node at (0,-2.0) {\small (b) augmented graph};
  \end{tikzpicture}
  \caption{The FJ model on the directed cycle of Figure~\ref{fig:networks}(a) when
  agent~$1$ is the only prejudiced agent ($\lambda_1 < 1$, $\lambda_i = 1$ for
  $i \neq 1$), and its augmented graph. Unlabelled edges have weight $1$. (a)~The
  prejudice $u_1$ enters agent~$1$ as an external input, with weight $1-\lambda_1$;
  this arrow is not an edge of $G$. (b)~The input becomes the virtual agent $1'$
  (dashed), which is stubborn and keeps the opinion $u_1$. The virtual agents
  $2', \dots, 6'$ are stubborn agents that nobody listens to, and are not drawn.}
  \label{fig:augmented}
\end{figure}
% REVISED (2026-10-24): Figure fig:augmented now has two panels, (a) the FJ model with
% the prejudice as an input arrow entering agent 1 and (b) the augmented graph
% (option (a) chosen by the user; the change of direction is explained in the text).

The augmented graph also shows, from the side of the FDG model, why consensus is not
the right notion for the FJ model. This observation is ours. The augmented
model~\eqref{eq:fj-augmented} is an FDG model, so the classification of
Table~\ref{tab:rosetta} applies to it. Each virtual agent is influenced only by
itself, so it is a sink of the condensation digraph of the augmented graph, and it is
aperiodic because of its self-loop. Hence the augmented graph has at least $n$ sinks
and it is not strongly connected; for $n \geq 2$, the consensus cases~(i) and~(iii) of
Table~\ref{tab:rosetta} never occur, and neither does case~(ii). Only two outcomes
remain.
\begin{itemize}
  \item \emph{Disagreement} (case~(iv)), when all the sinks are aperiodic. If all
        agents are P-dependent, the only sinks are the $n$ virtual agents, since every
        real agent is influenced, directly or through intermediaries, by some
        prejudiced agent $j$, and hence by $j'$. This is the situation of
        Remark~\ref{rem:many-stubborn}: by Theorem~\ref{thm:disagreement}, every real
        agent converges to a convex combination of the opinions of the stubborn
        virtual agents, i.e.\ of the prejudices, which is
        exactly~\eqref{eq:fj-V-all}. If some agents are P-independent, the sinks of
        $C^{22}$ are further sinks of the augmented graph, and all of them are
        aperiodic if and only if $W^{22}$ is regular.
  \item \emph{No convergence} (case~(v)), when at least one sink of $C^{22}$ is
        periodic, i.e.\ when $W^{22}$ is not regular.
\end{itemize}
These two outcomes are exactly the two alternatives of
Theorem~\ref{thm:fj-general}. The disagreement of case~(iv) concerns all the $2n$ agents
of the augmented model, whose virtual agents keep their different values $u_i$ forever
($x'(0) = u$ is part of the initial condition); it does not imply that the final
opinions of the real agents differ, which, as discussed above, depends on $V$ and $u$.
% REVISED (2026-10-28, review): shortened, same content. Was:
% These two outcomes are exactly the two alternatives of
% Theorem~\ref{thm:fj-general}. Note that the disagreement of case~(iv) concerns all the
% $2n$ agents of the augmented model: the virtual agents keep their different values
% $u_i$ forever, because in the augmented model the prejudices are part of the initial
% condition, $x'(0) = u$. It does not imply that the final opinions of the real agents
% differ, which, as discussed above, depends on $V$ and $u$.
% REVISED (2026-10-24): was "depends on u" (see the REVISED note above).
% NOTE (2026-10-24): footnote 21 of the tutorial (stubbornness as part of the state in
% the FDG model vs external input in the FJ model) goes to Section 3.2.1.

% TODO (to study before writing): the augmented graph~\cite{proskurnikov2017tutorial}:
% each prejudice u_i as a virtual stubborn agent anchored at u_i; TikZ figure; link with
% Remark rem:many-stubborn (several stubborn agents -> disagreement in the FDG model).
% TODO (OURS, R6, may be discarded): augmented graph = case (iv) of the Rosetta stone.
% TODO (to study with the simulations): consensus or disagreement. Sources: Bullo E5.24(v)
% (V_1 = [1 0; 1 0]: one prejudiced agent -> consensus; V_2: disagreement "for generic
% initial opinions"), tutorial Example 3 (cleavage). Consequences of V row-stochastic +
% zero columns (ours): one prejudiced agent, or equal prejudices of all prejudiced
% agents -> consensus. No general rule in the sources. Contrast with the FDG model
% (disagreement only from >= 2 sinks).
% DONE (2026-10-24): the TODOs above are written in the text, with the augmented graph,
% Figure fig:augmented, the consensus discussion (tutorial Definition 14) and our
% reading via Table 1 (R6). Excluded on purpose: Bullo E5.24(v) and the rule "one
% prejudiced agent <-> consensus" (not in the tutorial).

\subsubsection{Influence centrality}
\label{subsubsec:fj-centrality}

% TODO (to study before writing): influence centrality~\cite{proskurnikov2017tutorial}:
% asymptotically stable model and u = x(0), c = n^{-1} V^T 1_n (mean weight of x_i(0) in
% the final opinions), c^T 1_n = 1; Lemma 24 (Lambda = alpha I, alpha -> 1: c_alpha ->
% French's social power); possibly Bullo E5.24(iv) (Lambda = lambda I: A and V have the
% same left dominant eigenvector).
% DONE (2026-10-24): written below from Section 6.2 of the tutorial. Excluded on purpose:
% PageRank (not about social networks) and the reading of c on the augmented graph.

In the FDG model the social power $w$ is defined only when the group reaches a
consensus (Section~\ref{subsubsec:consensus}), which is not the relevant notion for the
FJ model (Section~\ref{subsubsec:fj-final}). Nor does the augmented FDG
model~\eqref{eq:fj-augmented} help: its graph has at least $n$ sinks, so for $n \geq 2$
it never reaches a consensus and its social power is not defined (this remark is ours). A different measure is needed,
and~\cite{proskurnikov2017tutorial} discusses an extension of the social power to the FJ
model, introduced by Friedkin.
% REVISED (2026-10-28, review): shortened; the n sinks of the augmented graph are shown in Section 3.1.3. Was:
% In the FDG model, the social power $w$ measures the weight of each initial opinion in
% the consensus value, and it is defined only when the group reaches a consensus
% (Section~\ref{subsubsec:consensus}). In the FJ model consensus is not the relevant
% notion (Section~\ref{subsubsec:fj-final}), so the social power cannot be used as it
% is. Rewriting the FJ model as the augmented FDG model~\eqref{eq:fj-augmented} does not
% help: the augmented graph has at least $n$ sinks, so for $n \geq 2$ the augmented
% model never reaches a consensus and its social power is not defined (this remark is ours). A
% different measure is needed:~\cite{proskurnikov2017tutorial} discusses an extension of
% the social power to the FJ model, introduced by Friedkin.

\paragraph{Influence centrality.}
As in~\cite{proskurnikov2017tutorial}, we confine ourselves to the asymptotically
stable FJ model with $u = x(0)$, the traditional assumption of the Friedkin--Johnsen
theory (Section~\ref{subsubsec:fj-model}); then $x(\infty) = V x(0)$
by~\eqref{eq:fj-V-all}. The assumption $u = x(0)$ is what gives the initial
opinions a role: in general $x(\infty) = Vu$ does not depend on $x(0)$, so the initial
opinions would have no weight in the final opinions. The social power of agent $i$ is
generalized as the mean weight of its initial opinion in determining the group
members' final opinions~\cite{proskurnikov2017tutorial}. These mean weights are the
elements of the non-negative vector $c$ such that
\begin{equation}\label{eq:influence}
  c = n^{-1} V^\top \mathbf{1}_n, \qquad
  \bar x = \frac{1}{n} \sum_{i=1}^{n} x_i(\infty)
         = \frac{1}{n} \mathbf{1}_n^\top V x(0) = c^\top x(0),
\end{equation}
and $c_i$ is called the \emph{influence centrality} of agent
$i$~\cite{proskurnikov2017tutorial}. In terms of $V$, $c_i$ is the average of the
$i$-th column of $V$, whose entries are the weights of $x_i(0)$ in the final opinions
of the single agents. Since $V$ is stochastic, $c$ is non-negative and $c^\top \mathbf{1}_n = 1$, like $w$
(Theorem~\ref{thm:consensus}), so the two measures can be compared directly. Similar to
French's social power, the influence centrality is generated by an opinion formation
mechanism~\cite{proskurnikov2017tutorial}.
% REVISED (2026-10-28, review): "sums to 1" was said twice. Was:
% Since $V$ is stochastic, $c^\top \mathbf{1}_n = 1$, and, similar
% to French's social power, the influence centrality is generated by an opinion
% formation mechanism~\cite{proskurnikov2017tutorial}. Like $w$ (Theorem~\ref{thm:consensus}),
% the vector $c$ is non-negative and sums to $1$, so the two measures can be compared
% directly.

The relation with the social power is made precise in~\cite{proskurnikov2017tutorial}.
The convergence and the consensus of the FDG model are equivalent, respectively, to
the regularity and the full regularity of the stochastic matrix
$A$~\cite{proskurnikov2017tutorial} (Definition~\ref{def:regular}). For a fixed $W$,
let $\Lambda = \alpha I_n$ with $\alpha \in (0,1)$, i.e.\ the same susceptibility for
all agents, and consider the corresponding matrix $V_\alpha$ and vector $c_\alpha$.
For $\alpha = 0$, i.e.\ $\Lambda = 0$, every agent keeps its initial opinion, $V = I$
and $c_0 = n^{-1} \mathbf{1}_n$: the social power is uniformly distributed between all
agents~\cite{proskurnikov2017tutorial}.

\begin{lemma}[\cite{proskurnikov2017tutorial}]
\label{lem:centrality-limit}
Let $W$ be a fully regular matrix and $w$ stand for the vector of French's social
power, corresponding to the FDG model~\eqref{eq:fdg} with $A = W$. Then,
$w = \lim_{\alpha \to 1-0} c_\alpha$.
\end{lemma}

The limit cannot be replaced by $\alpha = 1$: for $\Lambda = I$ the FJ model is the
FDG model, which is not asymptotically stable, and $I - W$ is singular since
$1 \in \operatorname{spec}(W)$ (Theorem~\ref{thm:stoch-spectrum}), so $V$ is not
defined. In social terms, the influence centrality moves from the uniform
distribution, when every agent is anchored to its own initial opinion, to French's
social power, when the weight of the prejudices vanishes.

% REMOVED (2026-10-23): \subsubsection{The general case: P-independent agents}
% \label{subsubsec:fj-general} -- merged into Section 3.1.2 (Theorem 21 and
% Corollary 23 of the tutorial), to follow the order of the tutorial.
